\needspace{9\baselineskip}
\begin{center}
\pdfbookmark[1]{3. Zu Art. 70.}{zusatz-003}
[3. Zu Art. 70.]
\end{center}
Da in den Gleichungen Art. 62 der Winkel \(N\) willkürlich ist, so setze man \(N=\lambda\); dadurch werden jene, wegen \(r'=r\cos\beta\), \(R'=R\cos B\), \(\Delta'=\Delta\cos b\):
\[
\begin{aligned}
1)\quad&r\cos\beta-R\cos B\cos(L-\lambda)=\Delta\cos b\cos(l-\lambda)\\
2)\quad&R\cos B\sin(\lambda-L)=\Delta\cos b\sin(l-\lambda)\\
3)\quad&r\sin\beta-R\sin B=\Delta\sin b.
\end{aligned}
\]

\SourceLeaf{302}Aus 2 und 3 folgt
\[
\operatorname{tang}(l-\lambda)=\frac{R\cos B\sin(\lambda-L)}{r\cos\beta-R\cos B\cos(\lambda-L)}\,\mathpunct{\vcenter{\hbox{\rule{.45em}{.35pt}}}}
\]
oder näherungsweise
\[
4)\quad l-\lambda=\frac{R\cos B\sin(\lambda-L)}{r\cos\beta}\mathpunct{\cdot}
\]

Ferner folgt aus 1 und 3
\[
\begin{aligned}
&R\cos B\sin\beta\cos(\lambda-L)-R\sin B\cos\beta\\
&\qquad=\Delta(\sin b\cos\beta-\cos b\sin\beta\cos(l-\lambda))
\end{aligned}
\]
oder
\[
\begin{aligned}
5)\quad&\frac{R\cos B\cos\beta(\operatorname{tang}\beta\cos(\lambda-L)-\operatorname{tang}B)}{r}\\
&\qquad=\frac{\Delta}{r}(\sin(b-\beta)+2\cos b\sin\beta\sin\tfrac12(l-\lambda)^2)
\end{aligned}
\]
oder näherungsweise
\[
=b-\beta.
\]
\begin{center}\rule{.18\textwidth}{.35pt}\end{center}
